\section{Introduction: the water that a crop uses}
A growing crop returns most of the water that reaches its root zone to the atmosphere, by \kw{transpiration} through its leaves and by \kw{evaporation} from the soil and from wet surfaces; the two together are \kw{evapotranspiration} (ET). ET is the largest outflow of the water balance of Chapter~2, it fixes the quantity of water that an irrigation system must supply, and it is the flux that a mulch or a trigger can change. It is driven by energy (the sun evaporates the water) and limited by supply (a dry soil cannot yield it). The standard method, developed by the FAO, separates these two aspects: the \emph{atmospheric demand} is expressed by a reference evapotranspiration $ET_0$ computed from weather alone; the \emph{crop} is expressed by a coefficient $K_c$ that scales $ET_0$ to its own demand $ET_c=K_cET_0$; and the \emph{stress} of a dry soil is a further factor $K_s\le1$ so that the actual rate is $ET_a=K_sK_cET_0$.

\section{Reference evapotranspiration: the Penman--Monteith equation}
The reference crop is a hypothetical grass of 0.12 m, with a fixed surface resistance of 70 s\,m$^{-1}$ and an albedo of 0.23, well watered and covering the ground. The FAO Penman--Monteith equation for a daily time step is

\begin{equation}ET_0=\frac{0.408\,\Delta\,(R_n-G)+\gamma\dfrac{900}{T+273}\,u_2\,(e_s-e_a)}{\Delta+\gamma\,(1+0.34\,u_2)},\label{eq:pm}\end{equation}


where $R_n$ is the net radiation (MJ\,m$^{-2}$\,d$^{-1}$), $G$ the soil heat flux (about zero daily), $T$ the mean temperature ($^\circ$C), $u_2$ the wind speed at 2 m (m\,s$^{-1}$), $e_s-e_a$ the vapour-pressure deficit (kPa), $\Delta$ the slope of the saturation vapour-pressure curve (kPa\,$^\circ$C$^{-1}$) and $\gamma$ the psychrometric constant (kPa\,$^\circ$C$^{-1}$). The first term of the numerator is the \emph{radiation} term, the second the \emph{aerodynamic} term; in humid climates the first dominates, in arid windy ones the second. Simpler equations (Hargreaves, Priestley--Taylor, pan evaporation with a pan coefficient) are used when data are missing; they are calibrated against Penman--Monteith.

\exbox{Penman--Monteith for a summer day}{Compute $ET_0$ for a summer day at 100 m altitude: $T_{max}=31$, $T_{min}=17$ $^\circ$C ($T=24$), $R_n=14.5$ MJ\,m$^{-2}$\,d$^{-1}$, $u_2=2.0$ m\,s$^{-1}$, mean relative humidity 55 \%.}\sol \givenline{$T_{max}=31$ ; $T_{min}=17$ $^\circ$C ; $T=24$ ; $R_n=14.5$ MJ\,m$^{-2}$\,d$^{-1}$ ; $u_2=2.0$ m\,s$^{-1}$}\ans Saturation vapour pressures: $e^\circ(31)=4.493$, $e^\circ(17)=1.938$, so $e_s=3.215$ kPa and $e_a=0.55e_s=1.768$ kPa; deficit $1.447$ kPa. $\Delta=4098\times0.6108\,e^{17.27\cdot24/(24+237.3)}/(24+237.3)^2=0.1791$. Pressure $P=101.3[(293-0.0065\times100)/293]^{5.26}=100.1$ kPa, $\gamma=0.000665P=0.0666$. Numerator: $0.408\times0.1791\times14.5=1.060$; $\gamma(900/297)\times2.0\times1.447=0.584$; sum $1.643$. Denominator: $0.1791+0.0666(1+0.68)=0.2910$. $ET_0=5.65$ mm\,d$^{-1}$ \ans The radiation term gives 3.6 mm and the aerodynamic term 2.0 mm of the total. \finans

\section{The crop coefficient and the four growth stages}
The crop coefficient $K_c$ changes through the season with the cover of the canopy. The season is divided into four stages: \emph{initial} (germination and early growth, length $L_1$, $K_c=K_{c,ini}$ nearly constant), \emph{development} ($L_2$, $K_c$ rises linearly to $K_{c,mid}$), \emph{mid-season} ($L_3$, $K_c=K_{c,mid}$) and \emph{late season} ($L_4$, $K_c$ falls linearly to $K_{c,end}$). The crop library of the application holds $K_{c,ini}$, $K_{c,mid}$, $K_{c,end}$, the four lengths, the maximum height $h$, the rooting depth $R_{max}$, the depletion fraction $p$ and the yield-response factor $K_y$ for 48 crops. the table below shows six of them.

\begin{table}[!htbp]\centering\small\caption{Six crops of the library: coefficients, stage lengths, height, rooting depth, depletion fraction and yield-response factor.}
\begin{tabular}{lrrrlrrrr}\toprule
\hd{Crop} & \hd{$K_{c,ini}$} & \hd{$K_{c,mid}$} & \hd{$K_{c,end}$} & \hd{Stage lengths (d)} & \hd{$h$ (m)} & \hd{$R_{max}$ (cm)} & \hd{$p$} & \hd{$K_y$}\\\midrule
Maize & 0.3 & 1.2 & 0.6 & 15+30+40+20 = 105 & 1.5 & 135 & 0.5 & 1.28\\
Wheat & 0.4 & 1.15 & 0.41 & 20+30+45+25 = 120 & 0.9 & 150 & 0.5 & 1.05\\
Tomato & 0.6 & 1.1 & 1.0 & 30+40+40+25 = 135 & 0.6 & 120 & 0.4 & 1.2\\
Potato & 0.5 & 1.1 & 0.4 & 25+30+45+30 = 130 & 0.6 & 60 & 0.4 & 1.1\\
Onion & 0.7 & 1.05 & 0.7 & 15+25+70+40 = 150 & 0.5 & 60 & 0.3 & 1.18\\
Soybean & 0.4 & 1.1 & 0.5 & 20+30+45+25 = 120 & 0.8 & 95 & 0.5 & 1.07\\\bottomrule\end{tabular}\end{table}

\exbox{Maize coefficient through the season}{For maize (the table below), compute $K_c$ on days 10, 30, 60, 90 and 110 after sowing and $ET_c$ if $ET_0=5.5$ mm\,d$^{-1}$ on each of those days.}\sol \givenline{$ET_0=5.5$ mm\,d$^{-1}$}\ans day 10: $K_c=0.300$, $ET_c=1.65$ mm; day 30: $K_c=0.750$, $ET_c=4.12$ mm; day 60: $K_c=1.200$, $ET_c=6.60$ mm; day 90: $K_c=1.050$, $ET_c=5.78$ mm; day 110: $K_c=0.600$, $ET_c=3.30$ mm \ans Day 10 is in the initial stage ($K_c=0.30$); day 30 in development (interpolated between 0.30 at day 15 and 1.20 at day 45: $0.30+0.90\times15/30=0.75$); 60 and 90 are in mid-season (1.20) and 110 is in late season: $1.20-0.60\times(110-85)/20$ is clipped at the end value 0.60 after day 105. \finans

\section{Evaporation from the soil: stages 1 and 2}
Soil evaporation follows two stages after a wetting (Ritchie, 1972). In \kw{stage 1} the surface is wet and the evaporation is limited by energy: it proceeds at the potential rate $E_p$ until the cumulative evaporation reaches $U$, the stage-1 capacity (about 5--12 mm, small in sand and large in clay). In \kw{stage 2} the surface layer has dried, evaporation is limited by the conductivity of the soil and falls as the square root of the time since stage 1 ended: $E(t)=\beta\,[\sqrt{t}-\sqrt{t-1}]$ per day, with $\beta$ a soil constant. The duration of stage 1 is $t_1=U/E_p$. A mulch lowers $E_p$ and hence lengthens stage 1 but also reduces the loss of stage 2; Chapter~10 treats it in detail. The two coefficients $K_e$ (evaporation) and $K_{cb}$ (basal, transpiration) of the dual-$K_c$ method split $K_c$ accordingly.

\exbox{Duration of stage 1}{The stage-1 capacity of a loam is $U=9$ mm and the potential evaporation is $E_p=4.5$ mm\,d$^{-1}$ in a hot week and 1.8 mm\,d$^{-1}$ in a cool one. Find the length of stage 1 in each, and the cumulative evaporation after 10 days in the hot week if $\beta=3.5$ mm\,d$^{-1/2}$.}\sol \givenline{$U=9$ mm ; $E_p=4.5$ mm\,d$^{-1}$ ; $\beta=3.5$ mm\,d}\ans $t_1=U/E_p$: hot week $9/4.5=2.0$ d; cool week $9/1.8=5.0$ d \ans After 10 days in the hot week: stage 1 gave 9 mm in 2 d; stage 2 lasts 8 d and gives $\beta\sqrt{8}=3.5\times2.83=9.9$ mm; total $=18.9$ mm \finans

\section{Water stress and yield response}
When the soil is drier than the allowed depletion the crop closes its stomata, transpires less than the demand and yields less. The FAO relation (Doorenbos and Kassam) relates the relative yield loss to the relative evapotranspiration deficit:

\begin{equation}1-\frac{Y_a}{Y_m}=K_y\Bigl(1-\frac{ET_a}{ET_c}\Bigr),\label{eq:ky}\end{equation}


with $K_y$ the yield-response factor: above 1 for crops that are sensitive to water (maize, 1.28; tomato, 1.20), below 1 for tolerant ones (sorghum, 0.90; groundnut, 0.70). The stress coefficient, within the daily model, is $K_s=(TAW-D_r)/[(1-p)TAW]$ once $D_r>p\,TAW=RAW$, and 1 before.

\exbox{Yield loss from a deficit}{A maize crop with $ET_c=520$ mm in the season gets $ET_a=430$ mm because irrigation was short. Find the relative yield and the yield if $Y_m=11.5$ Mg\,ha$^{-1}$. What is the yield of a tolerant crop with $K_y=0.7$ for the same deficit?}\sol \givenline{$ET_c=520$ mm ; $ET_a=430$ mm because irrigation was s ; $Y_m=11.5$ Mg\,ha$^{-1}$ ; $K_y=0.7$}\ans Relative deficit $1-430/520=0.173$. Maize: $1-Y_a/Y_m=1.28\times0.173=0.222$, so $Y_a/Y_m=0.778$ and $Y_a=8.95$ Mg\,ha$^{-1}$ \ans Tolerant crop: $1-0.7\times0.173=0.879$, so a loss of 12.1 per cent against 22.2 per cent for maize. \finans

